% ch-eu-fundamentals -- success labels on statutory fundamentals and the
% supervised prediction layer (imported from the handbook's success and
% gradient-boosting chapters; body of "Defining and Predicting Success
% from Fundamentals")

This chapter opens the transparent-regime track of the decision map (\cref{ch01:sec:map}). In
the European jurisdictions, statutory filings make firm fundamentals --- revenue, EBITDA,
payroll, balance-sheet items --- observable on a fiscal-year clock. The question ``is this firm
succeeding?'' therefore has a directly measurable answer and needs none of the point-process
machinery of the later parts: the target is a scalar success label, and the model that serves
it is supervised learning on a firm-year panel. Everything here is still downstream of
\cref{ch:data}: the deals table and macro feed supply features, and the point-in-time
discipline of \cref{prot:ch02:pit} governs every join. What changes is that a fundamentals
panel joins them --- and with it new ways to leak the future.

Against the defect ledger of \cref{tab:ch02:defects}, this regime trades defects rather than
escaping them. (C1) reappears as administrative censoring of recent cohorts' labels. (C3)
reappears as the survivor trap of \cref{cheu:sec:censoring}. (C5) is largely \emph{absent} ---
the statutory clock refreshes covariates whether or not the firm transacts, which is what makes
the regime transparent --- but a publication lag of 9--12 months takes its place, and
\cref{cheu:sec:worked} shows it is dangerous on its own. (C4) and (C6) do not bind: fiscal
years are trusted dates, and the filing universe is near-complete for limited-liability
companies.

\section{The fundamentals panel}\label{cheu:sec:panel}

For each EU firm $i$ and fiscal year $y$ we observe, from registry filings,
\begin{equation}\label{eq:success-panel}
F_{i,y} \;=\; \big(\text{revenue}_{i,y},\ \text{EBITDA}_{i,y},\
\text{payroll}_{i,y},\ \text{employees}_{i,y},\ \text{assets}_{i,y},\
\text{equity}_{i,y},\ \text{debt}_{i,y}\big),
\end{equation}
at yearly frequency, with near-complete coverage of limited-liability companies. The filing for
fiscal year $y$ is published at a date $\tau_{i,y}$ that lags fiscal year end by 9--12 months in
typical jurisdictions (statutory deadlines plus registry processing); $\tau_{i,y}$ is recorded
in the registry metadata and is itself firm-informative (late filers are disproportionately
distressed). Filings exist in \emph{vintages}: the as-reported figures published at
$\tau_{i,y}$, and later restatements that overwrite them in most commercial extracts.

\begin{warning}[Use as-reported vintages]\label{warn:success-vintage}
Restated figures incorporate information (audits, corrections, insolvency administrators'
revisions) that did not exist at $\tau_{i,y}$. Training features on restated fundamentals
violates the point-in-time discipline of \cref{prot:ch02:pit} exactly as consuming the
companies-table snapshot columns of \cref{ch02:sec:sources} would, and the bias is not neutral:
restatements are concentrated among firms that later failed, so restated features leak failure.
Features must be built from the as-reported vintage; restated data may be used for labels only
when the label is explicitly defined on final figures, and then the choice must be stated.
\end{warning}

\section{Candidate success metrics}\label{cheu:sec:metrics}

Fix a horizon $h$ (years) and write $E_y = \text{EBITDA}_{i,y}$, $R_y = \text{revenue}_{i,y}$,
$L_y = \text{employees}_{i,y}$, suppressing $i$.

\paragraph{EBITDA growth.} The naive candidate $\log(E_{y+h}/E_y)$ is undefined whenever $E_y
\le 0$ or $E_{y+h} \le 0$ --- for private growth companies, roughly a third of the sample.
Dropping those firm-years selects on profitability and destroys the target's meaning. We
instead use the \emph{arc (symmetric) growth rate}:

\begin{definition}[Arc growth]\label{def:success-arc}
For a fundamental $F$ with values $F_y$, $F_{y+h}$ not both zero,
\begin{equation}\label{eq:success-arc}
g^{(h)}_y \;=\; \frac{F_{y+h} - F_y}{\tfrac{1}{2}\big(\abs{F_{y+h}} +
\abs{F_y}\big)} \;\in\; [-2,\, 2],
\end{equation}
with $g^{(h)}_y = 0$ when $F_{y+h} = F_y = 0$.
\end{definition}

\begin{proposition}\label{prop:success-arc}
The arc growth rate is defined for all sign patterns, is bounded in $[-2,2]$, is antisymmetric
under swapping $F_y \leftrightarrow F_{y+h}$, and for $F_y, F_{y+h} > 0$ satisfies $g^{(h)}_y =
\log(F_{y+h}/F_y) + O\big(\log^3(F_{y+h}/F_y)\big)$, so it agrees with log-growth to second
order near zero growth.
\end{proposition}

\begin{proof}[Proof sketch]
Boundedness and antisymmetry are immediate from \cref{eq:success-arc}. For the expansion, set $u
= \log(F_{y+h}/F_y)$; then $g = 2(e^{u}-1)/(e^{u}+1) = 2\tanh(u/2) = u - u^3/12 + O(u^5)$. The
full argument, including the characterization of when $\abs{g}=2$ is attained, is in
\cref{app:proofs}.
\end{proof}

\begin{example}[Arc growth across a sign change]\label{ex:success-arc}
A firm improves from $E_y = -0.4$m to $E_{y+2} = +1.2$m. Log-growth is undefined; arc growth
gives $g = (1.2 - (-0.4)) / \tfrac{1}{2}(1.2 + 0.4) = 2.0$, the upper bound: sign-crossing
improvements saturate the scale, which is the intended reading (crossing into profitability is a
categorical success, and the metric does not pretend to rank magnitudes across the sign change).
By contrast a doubling from $0.6$m to $1.2$m gives $g = 0.6/0.9 = 0.667$ against $\log 2 =
0.693$: in the ordinary regime the two scales agree to within 4\%, as \cref{prop:success-arc}
predicts.
\end{example}

Boundedness is useful: a firm going from $E_y = 10$k to $E_{y+h} = 5$m is a success, but not
five hundred times the success of a doubling, and a bounded target prevents a few
small-denominator observations from dominating a squared-error fit. A
sign-preserving alternative when unbounded scale matters is the difference of
inverse-hyperbolic-sine transforms, $\operatorname{asinh}(F_{y+h}) - \operatorname{asinh}(F_y)$
with $\operatorname{asinh}(x) = \log(x + \sqrt{x^2+1})$, which behaves like the log-difference
for $\abs{F} \gg 1$ and remains defined through zero. In either case we winsorize the realized
target at its 1st and 99th percentiles, computed \emph{within training folds only}, to control
residual tail leverage.

\paragraph{Revenue CAGR.} $\mathrm{CAGR}^{(h)}_y = (R_{y+h}/R_y)^{1/h} - 1$, defined since $R_y
> 0$ for operating firms (pre-revenue firm-years are flagged and excluded from this metric,
itself a selection to be aware of). Robust to margin accounting; blind to profitability.

\paragraph{Margin expansion.} $\Delta m^{(h)}_y = E_{y+h}/R_{y+h} - E_y/R_y$, defined where both
revenues are positive. Captures operating-leverage improvement; extremely heavy-tailed for small
$R$ and sensitive to cost-capitalization choices.

\paragraph{Employment growth.} Arc growth \cref{eq:success-arc} applied to $L_y$. Nearly immune
to accounting discretion and currency; economically coarser.

\paragraph{Survival.} $S^{(h)}_y = \ind{\text{firm}~i~\text{is not insolvent or liquidated by
fiscal year}~y{+}h}$, from registry status events. Unambiguous and cheap, but nearly
uninformative in the bulk: most firms survive two years, so the indicator separates the left
tail only.

\begin{table}[htbp]
\centering
\small
\setlength{\tabcolsep}{4pt}
\begin{tabular}{lcccc}
\toprule
Metric & Interpretability & Acct.\ robustness & Tail behaviour & ML
target \\
\midrule
EBITDA arc growth      & high     & moderate & bounded        & high \\
EBITDA log-growth      & high     & moderate & undefined $E \le 0$ &
unusable \\
Revenue CAGR           & high     & high     & heavy right tail &
moderate \\
Margin expansion       & moderate & low      & very heavy      & low \\
Employment arc growth  & moderate & high     & bounded        & high \\
Survival indicator     & high     & high     & degenerate bulk & low
(imbalanced) \\
\bottomrule
\end{tabular}
\caption{Candidate success metrics. ``Acct.\ robustness'' is robustness
to discretionary accounting choices (depreciation schedules, cost
capitalization, revenue recognition); ``ML target'' summarizes
definedness, tail control, and informativeness for ranking.}
\label{cheu:tab:metrics}
\end{table}

\section{The formal estimand and its timing}\label{cheu:sec:estimand}

\begin{definition}[Success label]\label{def:success-label}
Fix a metric $g$ from \cref{cheu:sec:metrics} and horizon $h$. The label of firm $i$ at base
year $y$ is $Y_{i,y} = g\big(F_{i,y}, F_{i,y+h}\big)$, the metric evaluated over $(y, y+h]$. The
\emph{information date} of the pair is $\tau_{i,y}$, the publication date of the year-$y$
filing; a feature vector $Z_{i,\tau_{i,y}}$ is admissible if it satisfies \cref{prot:ch02:pit}
at $\tau_{i,y}$.
\end{definition}

The predictive estimand of the second half of this chapter is $\E[\,Y_{i,y} \st
Z_{i,\tau_{i,y}}\,]$; the causal estimand of \cref{ch:eucausal} replaces conditioning with
intervention on a financing exposure, $\E[\,Y_{i,y} \st \Do(A_{i,y}{=}a)\,]$. Both require $Y$
to be a single well-defined scalar attached to a single well-defined date \citep{hernan2020}.
Four label-engineering pitfalls recur in practice:

\begin{enumerate}
\item \textbf{Horizon misalignment.} Features dated at fiscal year end $y$ instead of at
$\tau_{i,y}$ grant the model 9--12 months of foresight; conversely a label read at
$\tau_{i,y+h}$ but indexed as ``$h$-year growth'' actually spans up to $h{+}1$ years of calendar
time. We index everything by $\tau$ and record the realized calendar span.
\item \textbf{Fiscal-year heterogeneity.} Fiscal year ends vary across firms and jurisdictions;
``year $y$'' is not a common calendar interval. Cross-sectional comparisons and time-based
validation splits must use $\tau_{i,y}$, not $y$.
\item \textbf{Consolidation changes.} A firm switching between solo and consolidated accounts
(or acquiring a subsidiary) exhibits spurious jumps in every fundamental. We flag
consolidation-scope changes from registry metadata and either exclude the spanning label or
difference within a constant scope.
\item \textbf{Currency.} Fundamentals are filed in local currency. Growth rates of the form
\cref{eq:success-arc} are currency-invariant only if both endpoints are in the same currency; we
compute growth in local currency and never convert at two different dates' exchange rates.
\end{enumerate}

\section{Censoring and the survivor trap}\label{cheu:sec:censoring}

A label $Y_{i,y}$ at horizon $h$ requires the year-$(y{+}h)$ filing. It is missing for two very
different reasons: the cohort is young (the calendar has not yet produced the filing ---
administrative censoring, the panel analogue of defect (C1)) or the firm exited (insolvency,
acquisition, migration out of filing obligations --- the analogue of (C3), here at least
partially \emph{labeled} by registry status). Simply dropping unlabelled firm-years conditions
the training distribution on surviving $h$ more years, which biases every estimate toward the
survivor population --- precisely the firms for which success is overrepresented, in the same
way that conditioning on the risk set must respect who is still observable \citep{cox1972}. The
risk-set machinery reappears for event data in \cref{ch:hawkes}; here a simple reweighting
suffices.

Let $C_{i,y} = \ind{F_{i,y+h}~\text{observed}}$ and let $\hat G(x_{i,y})$ estimate
$\Prob(C_{i,y}=1 \st x_{i,y})$ from pre-base-year covariates $x_{i,y}$ (age, size, sector,
filing punctuality), e.g.\ by logistic regression within calendar cohorts. The
inverse-probability-of-censoring (IPC) weighted loss
\begin{equation}\label{eq:success-ipcw}
\hat{L}(f) \;=\; \frac{1}{n}\sum_{(i,y)} \frac{C_{i,y}}{\hat
G(x_{i,y})}\, \ell\big(Y_{i,y},\, f(Z_{i,\tau_{i,y}})\big)
\end{equation}
is unbiased for the full-population risk when censoring is independent of $Y$ given $x_{i,y}$
--- an assumption to state, not to assume silently: exit by insolvency is plainly
outcome-dependent. Practically we (i) use \cref{eq:success-ipcw} for administrative censoring,
which is mechanically independent of $Y$ given the cohort; (ii) treat insolvency exits not as
censored but as \emph{observed failures}, assigning $Y = -2$ (the lower bound of arc growth)
with a sensitivity band; and (iii) treat acquisition exits, which are plausibly successes, as a
competing outcome reported separately rather than folded into $Y$. Truncating weights at $\hat G
\ge 0.1$ keeps the variance of \cref{eq:success-ipcw} controlled.

\section{Stylized facts to verify on any real panel}\label{cheu:sec:stylized}

Before fitting anything, three distributional facts should be confirmed on the panel at hand;
each has a design consequence, and their absence usually indicates a data problem rather than an
interesting economy.

\begin{enumerate}
\item \textbf{Heavy-tailed growth.} Firm growth rates are far from Gaussian: the cross-sectional
distribution of \cref{eq:success-arc} is sharply peaked, with tails well approximated by a
Laplace or heavier law. Consequence: squared error concentrates on the tails; prefer
pseudo-Huber or quantile losses and report rank metrics alongside RMSE (\cref{cheu:sec:why}).
\item \textbf{Mean reversion.} Successive growth rates correlate negatively: extreme years
partially reverse. Consequence: last year's growth is a strong but \emph{sign-flipping}
feature, and naive momentum readings of feature importances mislead; evaluation must be
out-of-time, not just out-of-sample.
\item \textbf{Size dependence.} Growth dispersion falls with size, roughly $\mathrm{sd}(g)
\propto R_y^{-\beta}$ with $\beta \approx 0.15$ (a well-documented deviation from Gibrat's
law). Consequence: any loss that ignores this heteroskedasticity overweights small firms;
stratify evaluation by size decile and consider variance-stabilizing weights.
\end{enumerate}

\Cref{cheu:tab:stylized} lists illustrative moments of the pattern these facts predict.
Verifying them on the synthetic EU panel (\cref{ch02:sec:synthetic}) is part of the validation
plan of \cref{ch:results-synth}; the same four columns should also be computed on any real
panel before the first model is fit.

\begin{table}[htbp]
\centering
\small
\begin{tabular}{lcccc}
\toprule
Metric ($h=2$) & Skewness & Excess kurtosis & Lag-1 autocorr. &
Undefined share \\
\midrule
EBITDA arc growth      & $-0.3$ & $9.8$  & $-0.21$ & $0.0\%$ \\
EBITDA log-growth      & ---    & ---    & ---     & $34\%$ \\
Revenue CAGR           & $+2.1$ & $14.6$ & $-0.12$ & $6\%$ \\
Margin expansion       & $-0.8$ & $27.3$ & $-0.35$ & $6\%$ \\
Employment arc growth  & $+0.4$ & $6.2$  & $-0.09$ & $0.0\%$ \\
\bottomrule
\end{tabular}
\caption{Illustrative stylized-fact moments (winsorization off, all
firm-years with both endpoints observed), to be verified on the
synthetic EU panel (\cref{ch:results-synth}). ``Undefined share''
is the fraction of firm-years on which the metric does not exist; the
log-growth row is blank because its moments are conditional on a
selected subsample and hence not comparable.}
\label{cheu:tab:stylized}
\end{table}

These facts jointly motivate the loss and cross-validation architecture of the prediction layer
below: heavy tails pick the loss, mean reversion picks out-of-time splits, size dependence picks
the stratification.

\section{Worked example: labelling a cohort}\label{cheu:sec:worked}

The definitions above compress into a short routine that is easy to get wrong. We run it once,
by hand, for a single firm whose EBITDA path crosses zero.
\Cref{cheu:tab:timeline} gives the firm's as-reported figures and registry publication dates;
fiscal years end on December 31.

\begin{table}[htbp]
\centering
\small
\begin{tabular}{lccc}
\toprule
Fiscal year $y$ & FY end & EBITDA $E_y$ (EUR m) & Published
$\tau_{i,y}$ \\
\midrule
2016 & 2016-12-31 & $-0.8$ & 2017-10-14 \\
2017 & 2017-12-31 & $-0.3$ & 2018-09-30 \\
2018 & 2018-12-31 & $+0.5$ & 2019-11-22 \\
2019 & 2019-12-31 & $+1.1$ & 2020-10-05 \\
2020 & 2020-12-31 & $+1.4$ & 2021-09-18 \\
\bottomrule
\end{tabular}
\caption{Timeline of the worked firm. Publication lags fiscal year end
by 9--11 months; the late 2018 filing (almost 11 months) is itself
informative (\cref{cheu:sec:panel}).}
\label{cheu:tab:timeline}
\end{table}

\paragraph{The label.} Take base year $y = 2017$ and the primary horizon $h = 2$. By
\cref{def:success-label} with arc growth \cref{eq:success-arc},
\begin{equation}\label{eq:success-worked-label}
Y_{i,2017} \;=\; \frac{E_{2019} - E_{2017}}{\tfrac{1}{2}\big(
\abs{E_{2019}} + \abs{E_{2017}}\big)}
\;=\; \frac{1.1 - (-0.3)}{\tfrac{1}{2}(1.1 + 0.3)} \;=\; 2.0,
\end{equation}
the saturated value: the firm crossed into profitability over $(2017, 2019]$, a categorical
success in the reading of \cref{ex:success-arc}.

\paragraph{Information-date alignment.} The pair is indexed at $\tau_{i,2017}$, i.e.\
2018-09-30, the instant the year-2017 figures become public. Admissible features there: the
as-reported 2016 and 2017 filings, deal rows with \code{deal\_date} before that date, and macro
values published before it. The label needs $E_{2019}$, public only on 2020-10-05: any training
set assembled earlier carries this pair as administratively censored
(\cref{cheu:sec:censoring}), and the realized calendar span --- information date to label
availability --- is 24.2 months, recorded alongside the nominal $h = 2$.

\begin{warning}[The not-yet-published filing]\label{warn:success-unpublished}
Consider the neighbouring pair, base year 2018. Dating its features at fiscal year end
2018-12-31 --- the natural join key in a panel indexed by fiscal year --- makes $E_{2018} =
+0.5$m a feature eleven months before its publication on 2019-11-22. The leaked content is
precisely the label's content: the sign crossing. A backtest learns that ``latest EBITDA just
turned positive'' predicts a saturated label; in production at the end of 2018 the newest public
figure is $E_{2017} = -0.3$m and the signal does not exist. The fix is mechanical, not
statistical: filings join as-of on $\tau_{i,y} \le t$, never on fiscal year --- the exact
analogue of joining macro series on \code{publish\_date} in \cref{prot:ch02:pit}(a) --- and the
staleness of the newest published filing enters as a feature, so the model sees lateness instead
of the future.
\end{warning}

The cohort routine is the by-hand computation industrialized: for every firm and every base year
whose filing is published, emit the pair at $\tau_{i,y}$; attach the label when the
year-$(y{+}h)$ filing publishes; score insolvency exits at $Y = -2$ and route acquisition exits
to the competing-outcome tally (\cref{cheu:sec:censoring}); winsorize realized labels within
training folds only. Every row of the resulting table carries $(i, y, \tau_{i,y}, Y, C)$, and
the validation splits of \cref{cheu:sec:validation} operate on $\tau$, never on $y$.

\subsection{Robustness of the label choice}\label{cheu:sec:robust}

The label is a construction, so before committing we ask how much rides on it.
\Cref{cheu:tab:robust} states the comparison to run on the synthetic panel
(\cref{ch:results-synth}), with the pattern we expect: the primary label against three
alternatives, on three stability criteria --- Spearman rank correlation of firm rankings with
the primary label; sensitivity of top-decile membership to moving the winsorization quantiles
from $1\%/99\%$ to $0.5\%/99.5\%$; and tail behaviour of the realized target.

\begin{table}[htbp]
\centering
\small
\setlength{\tabcolsep}{4pt}
\begin{tabular}{lccc}
\toprule
Label ($h=2$) & Rank corr. & Winsor.\ sensitivity & Tail behaviour \\
\midrule
EBITDA arc growth (primary) & $1.00$ & low      & bounded, mass at
$\pm 2$ \\
asinh-difference of EBITDA  & $0.86$ & moderate & unbounded, moderate
tails \\
Revenue CAGR                & $0.55$ & high     & heavy right tail \\
Employment arc growth       & $0.48$ & low      & bounded, thin tails \\
\bottomrule
\end{tabular}
\caption{Expected stability of the label choice, to be verified on the
synthetic panel (\cref{ch:results-synth}). ``Rank corr.'' is the
Spearman correlation of firm rankings with the primary label;
``winsor.\ sensitivity'' is the extent to which top-decile membership
changes when the winsorization quantiles move from $1\%/99\%$ to
$0.5\%/99.5\%$.}
\label{cheu:tab:robust}
\end{table}

The expected pattern is the same on any real panel: the asinh alternative reorders little (it
is a monotone reshaping of the same fundamental), while revenue- and employment-based labels
agree with the primary only moderately --- they measure different economics, not noisy copies
of one number. Rank correlations near $0.5$ mean that model comparisons, feature importances,
and top-decile readings can flip across definitions. We therefore re-run the headline models of
this chapter and \cref{ch:eucausal} under one pre-declared alternative --- employment arc
growth, chosen as the most accounting-immune and the most distant --- and report both. A
conclusion that survives both labels is about firms; one that does not is about the metric.
Declaring the single alternative before results are seen keeps this a robustness check rather
than metric shopping, in the same spirit as controlling a family of tests instead of reporting
the best one \citep{romano2019}.

The committed primary metric --- EBITDA arc growth at $h=2$, winsorized within folds, insolvency
exits scored at $-2$ --- is recorded with its rationale as decision D-EU.1 at the end of the
chapter. Three points of the rationale matter for the rest of the chapter. First, $h=2$
balances signal against censoring: at $h=1$ the label is dominated by accounting noise and mean
reversion (\cref{cheu:sec:stylized}); at $h=3$ the administratively censored share of recent
cohorts grows large and the IPC weights of \cref{eq:success-ipcw} get large. Second, EBITDA
rather than revenue keeps the metric sensitive to the margin dimension investors are paid on,
while the arc transform removes the sign-definedness problem that rules out log-growth. Third
--- and decisively for \cref{ch:eucausal} --- a composite index $Y^{\mathrm{comp}} = \sum_j w_j
Y_j$ degrades identifiability of the causal estimand: each component has its own confounding
structure, so a backdoor set valid for the composite need not exist even when one exists per
component; treatment effects of opposite sign on components cancel uninterpretably; and the
weights $w_j$ become untestable modelling choices inside the estimand itself. One outcome, one
horizon, one information date: everything downstream inherits its identifiability from this
choice.

\section{The learning problem}\label{cheu:sec:problem}

With the label fixed, the first modelling layer is plain supervised learning. A row is a
firm-year $(i,y)$. The label is $Y_{i,y} = w_c\!\bigl(g^{(h)}_{i,y}\bigr)$, the arc-growth of
EBITDA between statement years $y$ and $y+h$ with $h=2$, winsorized at level $c$ exactly as
above. The feature vector $X_{i,y} \in \R^p$ is assembled \emph{point-in-time}: it may use only
filings published and macro prints released by the information date $t_{\mathrm{info}}(i,y) =
\tau_{i,y}$ of the row, respecting \cref{prot:ch02:pit}. We estimate two objects:
\begin{equation}\label{eq:xgb-targets}
m(x) = \E[Y \st X = x],
\qquad
q_\tau(x) = \inf\{a : \Prob(Y \le a \st X = x) \ge \tau\},
\end{equation}
the conditional mean (point screening) and a small grid of conditional quantiles, $\tau \in
\{0.05, 0.10, 0.50, 0.90, 0.95\}$. The quantiles matter: the outcome is heavy-tailed even after
winsorization, upside screening acts on $q_{0.9}$ rather than the mean, downside risk control
acts on $q_{0.1}$, and the pair $(q_{0.05}, q_{0.95})$ is the input to the conformal intervals
of \cref{cheu:sec:uq}.

Two dependence structures dominate. \emph{Within-firm}: rows $(i,y)$ and $(i,y')$ share
persistent firm effects (management quality, niche economics), so errors are serially correlated
within $i$. \emph{Within-era}: rows of a given year $y$ share the macro shock, so errors are
cross-sectionally correlated within $y$. The first dictates firm-grouped resampling; the second
dictates era-blocked validation. Neither is optional, and both recur --- in bucketed form --- in
\cref{ch:buckets}.

\section{Feature catalogue}\label{cheu:sec:features}

\Cref{cheu:tab:features} lists the feature families. All are computable from the fundamentals
panel $F_{i,y}$, the deals table where coverage exists, and the macro feed; all obey the as-of
rule.

\begin{table}[htbp]
\centering
\caption{Feature catalogue for the EU prediction layer. All features
are as-of $t_{\mathrm{info}}(i,y)$; macro features use
publication-lagged vintages per \cref{prot:ch02:pit}(a).}
\label{cheu:tab:features}
\small
\begin{tabular}{@{}lll@{}}
\toprule
Family & Examples & Construction \\
\midrule
Lagged levels & EBITDA, revenue, payroll, assets & $\log(1+\cdot)$
where positive; \\
 & at $y, y{-}1, y{-}2$ & raw for signed EBITDA \\
Growth rates & 1--3y arc-growth of revenue, & same arc transform as
the \\
 & EBITDA, payroll & target, unwinsorized \\
Ratios & EBITDA margin, leverage, & clipped at 1st/99th pctile \\
 & asset turnover, payroll share & per sector-year \\
Peer position & within-sector-year rank and & computed per $(s,y)$
cell, \\
 & z-score of margin, growth, size & rank in $[0,1]$ \\
Sector-year aggregates & median growth, IQR, firm count & cell-level,
merged back \\
Firm statics & age at $y$, size bucket, country & age in years, log
assets bucket \\
Macro as-of & policy rate level and 12m change, & vintage at
$t_{\mathrm{info}}$; \\
 & CPI YoY, unemployment level/change & level \emph{and} change \\
Deal history & raised-to-date, round count, & zero-filled with
indicator \\
 & months since last round & where firm never in deals \\
\bottomrule
\end{tabular}
\end{table}

\paragraph{Missingness.} Fundamentals gaps are informative: small companies file abbreviated
accounts, distressed companies file late, and the propensity to be missing depends on the
(unobserved) value --- textbook MNAR. We therefore do \emph{not} mean-impute: imputation
overwrites the signal carried by absence and biases every downstream split. Instead we rely on
the native sparsity handling of gradient boosting, in which each split learns a default
direction for missing values from the data \citep{chenguestrin2016}, and we add explicit
indicators (\code{ebitda\_missing\_lag1}, \code{never\_in\_deals}) for gaps whose economic
meaning we can state. Trees can then interact the indicator with size and sector: ``missing
EBITDA at small size'' differs from ``missing EBITDA at large size'', and the model should be
free to learn that.

\section{Why boosted trees here --- and what they cannot do}\label{cheu:sec:why}

The learner is gradient-boosted trees in the XGBoost formulation
\citep{chenguestrin2016,friedman2001}; the second-order boosting step, gain formula, and
sparsity-aware default directions are derived once in \cref{ch:buckets} (\cref{prop:ch09:xgb})
and are not repeated here. The case for boosted trees on \emph{this} problem has four parts:
(i) the signal is riddled with interactions (margin $\times$ sector, growth $\times$ rate
regime) that we cannot enumerate a priori and trees discover automatically; (ii) the features
are mixed continuous/ordinal/categorical with wildly different scales, and trees are invariant
to monotone transforms of each feature, so we spend no effort on marginal calibration; (iii)
native missingness handling, as above; (iv) strong, well-understood regularization (depth,
shrinkage, subsampling, $\ell_2$ on leaf weights).

The case against has two entries, and both matter here. First, \emph{no extrapolation}: a tree
ensemble is piecewise constant, so its prediction outside the convex hull of the training
features is the value at the nearest boundary cell. If every training era has falling policy
rates, the fitted rate response is a flat extension --- the model cannot form an opinion about
a rising-rate regime it has never seen. Second, \emph{unstable attribution}: correlated
features split importance unpredictably across refits, so feature rankings are not portable
conclusions.

We mitigate rather than solve. Era-aware validation (\cref{cheu:sec:validation}) at least
measures degradation across regimes instead of hiding it. Monotone constraints encode signs we
are willing to defend economically --- e.g.\ growth non-increasing in the 12-month rate
\emph{change}, downside quantile non-increasing in leverage --- which reduces variance in thin
regions and prevents sign flips there.

\begin{remark}[Constraints are priors, not facts]\label{rem:xgb-mono}
A monotone constraint changes the model class, not the world. Impose one only where the sign
survives scrutiny (the rate \emph{level} is ambiguous: it proxies both financing cost and the
demand cycle; the \emph{change} is the cleaner candidate). A wrong constraint is worse than
none, because it silently converts local noise into global bias. The same discipline governs
the bucketed panel of \cref{ch:buckets}, where an apparently obvious monotonicity in staleness
age is deliberately declined.
\end{remark}

For the conditional-mean model, heavy tails argue against squared error even after
winsorization. We use the pseudo-Huber objective
\begin{equation}\label{eq:xgb-phuber}
\ell_\delta(r) = \delta^2\left(\sqrt{1 + (r/\delta)^2} - 1\right),
\end{equation}
which is quadratic for $\abs{r} \ll \delta$ and linear for $\abs{r} \gg \delta$, with $\delta$
set to roughly one cross-sectional standard deviation of $Y$. Note the fitted functional then
sits between conditional mean and median; we report which loss produced which model and never
mix them in evaluation.

\section{Validation design}\label{cheu:sec:validation}

The general temporal-evaluation constitution of this document --- fixed development split,
rolling-origin headlines, no post-split information anywhere --- is stated once in
\cref{ch07:sec:temporal} (\cref{prot:ch07:temporal}) and applies here unchanged. What is
EU-specific is the shape the protocol takes on a firm-year panel with overlapping two-year
labels, and why anything weaker fails.

\paragraph{Never random K-fold.} Random row-level K-fold fails three ways on this panel.
(i)~\emph{Firm leakage}: rows of the same firm land on both sides of the split, and the
persistent firm effect makes validation loss optimistic --- the model is partly memorizing
firms. (ii)~\emph{Target overlap}: with $h=2$, the labels of years $y$ and $y+1$ share a
statement year, so adjacent rows carry overlapping information about each other's targets.
(iii)~\emph{Era leakage}: the model trains on the future macro regime it will be asked to
predict, which no deployed model enjoys. All three inflate measured skill precisely where the
deployment risk is largest.

\paragraph{Rolling-origin splits with an embargo.} We order eras (label years) $y_1 < \dots <
y_E$ by information date and build folds $k = 1, \dots, K$: train on eras up to $y_{e_k}$,
\emph{embargo} the next $h$ eras (their labels overlap the training window), validate on the era
after the embargo. Every fold mimics deployment: fit on the past, predict a disjoint future. The
same firm may appear in train (early years) and validation (later years) --- that is legitimate
and realistic; the leakage to prevent is label overlap, handled by the embargo, not firm
recurrence across time. The embargo is the one ingredient the general protocol of
\cref{ch07:sec:temporal} does not need to state, because only horizon labels overlap this way.

\paragraph{Hyperparameter protocol.} Random search (100--200 draws) over the ranges in
\cref{cheu:tab:hparams}; each draw is scored by its \emph{average validation loss across all
rolling folds}, with the number of boosting rounds chosen by early stopping on the \emph{most
recent} fold only --- the one knob that benefits most from recency. The final model refits
on all data with the tuned configuration and the early-stopped round count scaled by the ratio
of full to training rows. Tuning on an average across eras deliberately trades a little peak
skill for regime stability.

\begin{table}[htbp]
\centering
\caption{Search ranges. Depth is kept shallow: firm-year panels reward
many weak interactions over few deep ones.}
\label{cheu:tab:hparams}
\begin{tabular}{@{}lll@{}}
\toprule
Parameter & Range & Note \\
\midrule
max depth & 3--6 & log-uniform integer \\
learning rate & 0.01--0.1 & log-uniform \\
min child weight & 5--100 & guards thin sector cells \\
row subsample & 0.6--0.9 & per boosting round \\
column subsample & 0.5--0.9 & decorrelates trees \\
$\ell_2$ leaf penalty & 0.5--20 & log-uniform \\
rounds & early stopping & patience 50, most recent fold \\
\bottomrule
\end{tabular}
\end{table}

\section{Uncertainty quantification}\label{cheu:sec:uq}

Three tools, three targets: per-firm heteroskedastic intervals (quantile regression),
finite-sample coverage repair (conformal), and model-level uncertainty of scalar summaries
(cluster bootstrap). \Cref{cheu:tab:uqtools} states when each is the right instrument.

\subsection{Quantile regression with the pinball loss}

For $\tau \in (0,1)$ the pinball (check) loss is
\begin{equation}\label{eq:xgb-pinball}
\rho_\tau(u) = u\,\bigl(\tau - \ind{u < 0}\bigr)
= \tau\, u_+ + (1-\tau)\, u_-,
\end{equation}
with $u_+ = \max(u,0)$, $u_- = \max(-u,0)$. Its population minimizer is exactly the quantile we
want:

\begin{proposition}[The pinball minimizer is the $\tau$-quantile]
\label{prop:xgb-pinball} Let $Y$ have distribution function $F$ and finite first moment, and let
$\varphi(a) = \E[\rho_\tau(Y - a)]$. Then any $a^*$ with $F(a^{*-}) \le \tau \le F(a^*)$
minimizes $\varphi$; if $F$ is continuous and strictly increasing, $a^* = F^{-1}(\tau)$
uniquely.
\end{proposition}

\begin{proof}
$\varphi$ is convex, and for $a$ a continuity point of $F$, $\varphi'(a) = -\tau \Prob(Y > a) +
(1-\tau)\Prob(Y < a) = F(a) - \tau$. A convex function is minimized where its (sub)derivative
changes sign, i.e.\ where $F(a^{-}) - \tau \le 0 \le F(a) - \tau$. The version with all
measure-theoretic care (general $F$, subgradient characterization) is in \cref{app:proofs}.
\end{proof}

We fit one boosted model per $\tau$ with objective \cref{eq:xgb-pinball} applied conditionally.
Separately fitted quantiles can cross in finite samples; we repair by sorting the predicted grid
per row (monotone rearrangement), which never worsens pinball loss.

\subsection{Conformalized quantile regression}

Quantile models are calibrated only insofar as the model class and sample size allow; conformal
prediction repairs marginal coverage with finite-sample guarantees \citep{romano2019}.

\paragraph{Algorithm (split CQR).}
\begin{enumerate}
\item Split the available rows into a proper training set $I_1$ and a calibration set $I_2$ of
size $n_2$, with the split at the firm level and $I_2$ chosen as a \emph{block}: the most recent
complete era(s) before the deployment date.
\item On $I_1$, fit lower and upper quantile models $\hat q_{\alpha/2}$, $\hat q_{1-\alpha/2}$
as above.
\item For each $j \in I_2$ compute the conformity score $E_j = \max\bigl\{\hat q_{\alpha/2}(X_j)
- Y_j,\; Y_j - \hat q_{1-\alpha/2}(X_j)\bigr\}$.
\item Let $Q_{1-\alpha}$ be the $\lceil (n_2+1)(1-\alpha) \rceil$-th smallest of $\{E_j\}_{j \in
I_2}$.
\item Report the interval $C(x) = \bigl[\hat q_{\alpha/2}(x) - Q_{1-\alpha},\; \hat
q_{1-\alpha/2}(x) + Q_{1-\alpha}\bigr]$.
\end{enumerate}

\begin{theorem}[CQR marginal coverage, \citealp{romano2019}]
\label{thm:xgb-cqr} If the calibration pairs $\{(X_j, Y_j)\}_{j \in I_2}$ and the test pair $(X,
Y)$ are exchangeable, then $\Prob\{Y \in C(X)\} \ge 1 - \alpha$; if the scores are almost surely
distinct, also $\Prob\{Y \in C(X)\} \le 1 - \alpha + 1/(n_2 + 1)$.
\end{theorem}

The proof, via an exchangeable-rank lemma, is in \cref{app:proofs}.

\begin{warning}[Temporal data are not exchangeable]\label{warn:xgb-exch}
\Cref{thm:xgb-cqr} assumes the test row is exchangeable with the calibration rows. A macro-era
shift between calibration block and deployment breaks this, and no finite-sample guarantee
survives. Practice: calibrate on the most recent contiguous block (as above), \emph{treat the
nominal level as approximate}, measure empirical coverage per era on every historical fold, and
recalibrate at every model vintage. If per-era coverage drifts systematically with the macro
state, widen intervals by the worst observed era's deficit rather than trusting the nominal
$\alpha$.
\end{warning}

\subsection{Cluster bootstrap over firms}

Quantile and conformal intervals describe \emph{per-row} outcome uncertainty. A different
question --- how uncertain is a scalar summary of the fitted model, such as an era MAE, a
decile-spread statistic, or the skill gap over a benchmark --- requires resampling the fitting
process itself \citep{efron1994}.

\paragraph{Algorithm (firm-level percentile bootstrap).}
\begin{enumerate}
\item Let $\firms$ be the set of firms. For $b = 1, \dots, B$ (we use $B \in [200, 500]$): draw
$\abs{\firms}$ firms from $\firms$ with replacement and assemble the replicate panel from
\emph{all rows} of each drawn firm (duplicates kept).
\item Re-run the frozen pipeline on the replicate: feature construction, the tuned
hyperparameter configuration (not re-searched), fit, and compute the scalar summary
$\hat\theta_b$ on the replicate's validation eras.
\item Report the percentile interval $\bigl[\hat\theta_{(\lceil \alpha B/2 \rceil)},\;
\hat\theta_{(\lceil (1-\alpha/2) B \rceil)}\bigr]$.
\end{enumerate}
Resampling firms, not rows, is the key choice: within-firm correlation means the effective
sample size is closer to the number of firms than the number of rows, and a row-level bootstrap
understates variance accordingly.

\begin{table}[htbp]
\centering
\caption{Choosing the uncertainty tool.}
\label{cheu:tab:uqtools}
\begin{tabular}{@{}llll@{}}
\toprule
Tool & Quantifies & Guarantee & Cost \\
\midrule
Quantile regression & per-row outcome spread & none per se & 1 fit per
$\tau$ \\
Split CQR & per-row interval coverage & finite-sample marginal, & 1
calibration \\
 & & approximate under drift & pass \\
Firm bootstrap & model-level scalars & asymptotic, within & $B$
refits \\
 & (MAE, skill gaps) & the observed design & \\
\bottomrule
\end{tabular}
\end{table}

\section{Evaluation protocol and benchmarks}\label{cheu:sec:eval}

We report MAE and $R^2$ \emph{by era} and by horizon, never pooled alone: a pooled MAE averages
over exactly the regime heterogeneity the model is weakest on. The median model (pinball at
$\tau = 0.5$) is evaluated with MAE --- matched loss --- and the pseudo-Huber mean model with
both MAE and RMSE. Quantile calibration is checked per era with reliability diagrams: for each
$\tau$, the empirical fraction of outcomes below $\hat q_\tau(X)$ plotted against $\tau$; a
well-calibrated model tracks the diagonal, a regime-sensitive one shows era-dependent slope.
Subgroup stability: per-cell MAE by sector and size tercile relative to pooled MAE, flagging
cells where a benchmark beats the boosted model --- such cells are candidates for exclusion from
deployment, not for silent averaging. This is the panel-regression face of the reporting
discipline of \cref{ch07:sec:reporting}: no number without the era, subgroup, and benchmark that
qualify it.

\begin{warning}[Attributions are associational]\label{warn:xgb-shap}
SHAP-style attributions decompose the \emph{fitted} $\hat m$, i.e.\ they redistribute predicted
differences among correlated inputs of an associational object. A large attribution on
\code{months\_since\_last\_round} does not mean financing drives growth; it means firms that
recently raised look different, which is precisely the selection effect \cref{ch:eucausal} is
built to separate from the causal one. Use attributions for debugging (leakage detection, sanity
of signs) and never in an investment memo as an effect size.
\end{warning}

Every fold carries three benchmarks: \emph{persistence} (predict the firm's last observed
growth), \emph{sector-year mean} (predict the cell average), and \emph{ridge} regression on the
identical feature matrix (with the missingness indicators, since ridge has no native sparsity
handling). The shipping rule is strict: the boosted model is deployed only where its out-of-era
skill gap over ridge exceeds the half-width of the firm-bootstrap interval on that gap. A gap
inside the interval is noise, and ridge is cheaper, more stable, and partially extrapolates.
\Cref{cheu:tab:notxgb} records the situations in which gradient boosting is the wrong tool
altogether.

\begin{table}[htbp]
\centering
\caption{When gradient boosting is the wrong tool.}
\label{cheu:tab:notxgb}
\begin{tabular}{@{}lll@{}}
\toprule
Situation & Symptom & Use instead \\
\midrule
Tiny sample & $\lesssim 5$k firm-years & ridge on curated features \\
Regime extrapolation & deployment macro outside & linear/constrained
macro \\
 & training support & term $+$ trees on the rest \\
Regulatory interpretability & auditable decision rule required &
scorecard or shallow \\
 & & monotone model \\
Skill gap within noise & bootstrap CI covers zero & ridge; revisit
features \\
\bottomrule
\end{tabular}
\end{table}

The model of this chapter answers ``which firms will grow?'' conditional on everything
observable. It is silent on ``what happens if a firm raises?'' --- the fitted dependence on
deal-history features mixes selection with effect. \Cref{ch:eucausal} takes up that question
with the same data and an explicitly causal contract.

\begin{decision}{D-EU.1 --- the transparent-regime target and its gate}
The primary success metric of the transparent-regime track is \emph{EBITDA arc growth at horizon
$h=2$} (\cref{def:success-arc}), winsorized at the within-fold 1\%/99\% quantiles, with
insolvency exits scored at $-2$, acquisition exits reported as a competing outcome, and
administrative censoring handled by the IPC weights of \cref{eq:success-ipcw}; no composite
index is used, and one pre-declared alternative label (employment arc growth) is re-run for
every headline result. The committed learner is gradient-boosted trees (pseudo-Huber mean plus a
pinball quantile grid) under the embargoed rolling-origin validation of
\cref{cheu:sec:validation}, and the deployment gate is the ridge benchmark: the boosted model
ships only in cells where its out-of-era skill gap over ridge exceeds the firm-bootstrap
half-width on that gap. Rationale: signal-vs-censoring balance at $h=2$, sign-definedness of the
arc transform, and single-outcome identifiability for \cref{ch:eucausal}; \cref{cheu:tab:robust}
for why the alternative label is mandatory. Reversal trigger: if the pre-declared alternative
label flips the sign or ordering of a headline conclusion in two consecutive evaluation
campaigns, the label choice is reopened and this box is rewritten.
\end{decision}

\begin{codehooks}
The cohort routine of \cref{cheu:sec:worked} becomes a planned panel builder
\modrepo{eu/panel.py}, emitting rows $(i, y, \tau_{i,y}, Y, C)$ with as-reported vintages only,
mirroring the loader contract of \code{synthetic/loaders.py} (as-of joins, staleness ages
carried as features). Label and metric implementations (arc growth, asinh difference,
within-fold winsorization) belong in a planned \modrepo{eu/labels.py} with the stylized-fact
audit of \cref{cheu:sec:stylized} as its test suite; the synthetic EU-style panel behind
\cref{cheu:tab:stylized,cheu:tab:robust} will be generated from the synthetic world in
\code{handbook/code/synthetic/} (\code{firms.py}, \code{marks.py}, \code{observe.py}). The
boosted stack --- embargoed rolling-origin folds, pinball grid, split CQR, firm-level bootstrap
--- is shared infrastructure with the bucketed panel of \cref{ch:buckets} and should live once,
in the planned \modrepo{panel/} learner package, not twice.
\end{codehooks}
